% ============================================================
% ch01 总览：论文全景图
% ============================================================
\section{总览：论文全景图}\label{sec:ch01}

\begin{keybox}{先抓住本章的四件事}
\begin{itemize}
  \item 论文检验的承诺是``尖锐度若可廉价测量，则可廉价操纵''——
        结论是这句承诺\textbf{一分为二}：测量端成立、操纵端结构性失败。
  \item 组织全文的轴是\textbf{预测器 $\neq$ 执行器}：同一个代理量，
        拿来监测 $\lmax$ 近乎精确，拿去驱动控制全部失败。
  \item 控制端失败的形态有三种：帽原地不动（钉在 $2/\eta$）、帽被
        顶开（上逃逸）、有动作仍被钉住（规模实验）——共同点是
        \textbf{永远压不到阈值以下}；唯有 SAM 式全局动作能下行。
  \item 论文的证据链是``实验 + 机制排除 + 2-D 约定理''的三层结构，
        每一层都有明确的适用域声明（``2-D 约定理''即二维约化定理：
        把百万维训练动力学压缩成两个变量再证明的定理，见第九章）。
\end{itemize}
\end{keybox}

\subsection{从教程到论文：问题的自然延伸}\label{sec:ch01-question}

教程收尾时你已经知道：全批梯度下降训练神经网络时，海森矩阵最大
特征值 $\lmax$ 会自发攀升到稳定阈值 $2/\eta$ 附近并振荡驻留
（Cohen 等人 2021 的 EoS 现象）。驱动它攀升的是\textbf{渐进尖锐化}
（progressive sharpening——训练中 $\lmax$ 自发增长的趋势，
Li 等人 2022）；把它按在阈值上的，则是\textbf{周期-2 振荡}
（period-2 oscillation——参数沿顶部方向每两步一来回的跳动）所带来
的\textbf{自稳定化}（self-stabilization——跳动的三阶泰勒项产生
恢复力、把 $\lmax$ 拉回 $2/\eta$，Damian 等人 2023）。教程也展示了
两个看似``可利用''的伴生现象：
训练轨迹是低维的（low-dimensional trajectory：PCA 第一主成分
解释 84--99\% 方差），而且权重范数等廉价
量似乎与 $\lmax$ 相关（Li 等人 2022 的代理）。

\textbf{术语提醒。}以上概念在本解读中直接使用；每个术语的英文原名与
一句话解释集中列在前言（第~\ref{sec:ch00}~章）的``术语与记号约定''与``背景速查''，
后文首次引入的新概念则在出现处就地标注。

把这些放在一起，一个诱人的想法几乎自动浮现：\textbf{既然尖锐度
可测、轨迹低维、还有廉价代理，我们能不能造一个优化器，实时监测
尖锐度并在它过高时把它压回去？}如果成功，EoS 就不只是理论好奇，
而是通向``平坦最小值（flat minima）$\Rightarrow$ 好泛化''这一经典信条
（Keskar 等人 2017）的工程抓手。SAM（Foret 等人 2021）的成功——
通过对梯度做全局扰动确实能训练到更平的最小值——让这条路线显得
更加可信。

\begin{notebox}[title={概念详解：SAM 是什么}]
SAM（\textbf{S}harpness-\textbf{A}ware \textbf{M}inimization，锐度感知
最小化；Foret 等人 2021）是一个为``找平坦最小值''而设计的优化器，
动机正是上面的平坦度--泛化信条。它的想法是把普通的``最小化当前点的
损失 $L(\theta)$''升级为``最小化当前点周围半径 $\rho$ 的邻域内的
\emph{最坏}损失''$\max_{\|\varepsilon\|\le\rho}L(\theta+\varepsilon)$
：不满足于脚下这一点的损失低，还要求旁边一圈都低——落点必须平坦。
直接解内层的极大化不可行，SAM 用一步一阶近似取极值方向：邻域内让
损失上升最快的方向就是梯度方向本身，于是取
\[
  \varepsilon^{*}=\rho\,\frac{\mathbf{g}}{\|\mathbf{g}\|},
  \qquad \mathbf{g}=\nabla L(\theta),
\]
即把参数沿\emph{归一化后的梯度}探出固定长度 $\rho$（$\rho$ 是 SAM
的邻域半径超参，本文取 $0.05$）。每步训练由此变成四拍：
\textbf{（1）}在 $\theta$ 算梯度 $\mathbf{g}$（第一次前向 + 反传）；
\textbf{（2）}构造扰动 $\varepsilon^{*}$（一次归一化，几乎免费）；
\textbf{（3）}在扰动点 $\theta+\varepsilon^{*}$ 再算梯度
$\nabla L(\theta+\varepsilon^{*})$（第二次前向 + 反传）；
\textbf{（4）}用\emph{这个}梯度对原参数 $\theta$ 做普通下降更新。
代价约为普通训练的两倍。

本文反复使用的``\textbf{全局动作}''一词，指的就是这个算法的数学实质
：把 $\nabla L(\theta+\varepsilon^{*})$ 做一阶展开（第~\ref{sec:ch09}~章
推论的形式化），得有效梯度
$\approx\bigl(I+\tfrac{\rho}{\|\mathbf{g}\|}H\bigr)\mathbf{g}$——海森
的\emph{每一个}特征方向上的分量都被乘上各自的因子
$1+\rho\lambda_i/\|\mathbf{g}\|$：曲率越大的方向被放大得越多。SAM 的
每一步都在\emph{全部}参数方向上同时做``按曲率加权''的重缩放；这与
秩-1 干预（只改写\emph{一个}方向上的梯度分量）构成了论文控制端的
整条对照轴。注意本文对 SAM 的用法是``控制端的逃逸基准''——唯一能把
$\lmax$ 压到阈值之下的对照臂——而不评价其泛化主张；SAM 家族变体
（ASAM、SAF 等）的范围限定见第~\ref{sec:ch02}~章与第~\ref{sec:ch10}~章。
\end{notebox}

论文正是对这一承诺的直接检验，而且开头就把检验结果摆明了
（论文摘要）：

\begin{quote}\itshape
``廉价的尖锐度\textbf{预测器}在 plain、BatchNorm、LayerNorm 架构
上都能可靠追踪顶部特征值……廉价的尖锐度\textbf{控制}则结构性失败：
沿顶部特征向量阻尼梯度分量的秩-1 修正步无法把 $\lmax$ 压到阈值
以下——它要么让自稳定的 $2/\eta$ 均衡原样保持（尖锐度钉在那里），
要么使均衡失稳（尖锐度反而冲过阈值）……而 SAM 的全局扰动确实能把
均衡压下去。''
\end{quote}

\noindent 引文里有两个词值得当场展开。\textbf{其一是 plain。}``plain
网络''指\textbf{不带任何归一化层}的网络——本文的 tanh MLP 就是 plain
网络。论文之所以按``有没有归一化、是哪一种''给架构分族（plain /
BatchNorm / LayerNorm），是因为归一化层改变了``网络的输出尺度由谁
控制''这件事：plain 网络里权重直接决定激活幅度，$\|\mathbf{W}\|$ 大
$\Rightarrow$ 激活大 $\Rightarrow$ 曲率大，权重范数与 $\lmax$ 天然
耦合——Li 等 (2022) 的范数代理正是在这种网络上提出并成立的。
BatchNorm（卷积网络标配）与 LayerNorm（Transformer 标配）则都在
``先把激活标准化、再施加可学习仿射''的结构下把这条传导链剪断，尺度
信息转移到自己的仿射参数（BN 的 $\gamma$）身上——机制与后果在第四章
（\ref{sec:ch04-tier0} 节）展开。因此``三族齐备''不是架构多样性的
装饰，而是 Tier-0 廉价代理可靠性地图的三块地质板块。术语陷阱：ResNet
文献里的 plain network（He 等 2016，``无残差连接''的对照网络）与
此处用法\emph{无关}——本文的 plain 只关乎归一化层。\textbf{其二是
均衡。}``均衡''指 EoS 轨道长期驻留的状态，见
第~\ref{sec:ch00-glossary}~节 D 组。

\subsection{预测端：一条代理阶梯}\label{sec:ch01-predictors}

论文的第一条腿（第四章解读）建立了一个``$\lmax$ 代理阶梯''，
按计算成本分层：

\begin{itemize}
  \item \textbf{Tier-0（零 HVP）}：各层权重范数、BN 的 $\gamma$ 范数
        （$\gamma$ 即 BN 层的逐通道缩放参数，见第~\ref{sec:ch00}~章
        术语表与第四章）、梯度范数。它们的可靠性\textbf{依赖归一化
        类型且跨种子不稳定}：plain 网络上权重范数可用（去趋势
        Pearson——先剔除共同漂移再算的相关，见
        第~\ref{sec:ch03-stats}~节——$0.65$--$0.71$，
        最后一个隐层最优）；BatchNorm 下权重范数相关由正转负（$-0.8$）、
        信号转移到 $\gamma$，但 $\gamma$ 多种子下去趋势相关
        $+0.247\pm0.340$、置信区间跨零；LayerNorm 的 Transformer 上
        所有零计算代理都触顶于 $\|\mathbf{g}\|$ 的 $+0.52$。
  \item \textbf{Tier-1（单 HVP）}：梯度瑞利商
        $\lhatg=\mathbf{g}^{\!\top}\!H\mathbf{g}/\mathbf{g}^{\!\top}\!\mathbf{g}$。
        在 ResNet-18 上 EoS 体制（$|\cos(\mathbf{g},\vtope)|>0.7$）的
        去趋势相关达 $+0.981\pm0.030$——\textbf{近乎精确}。这是论文
        交付给实践者的监测器。
  \item \textbf{顶端（约 5 HVP）}：热启动幂迭代（Ramasinghe 等人
        2026），作为 ground-truth 测量，论文不将其据为贡献。
\end{itemize}

预测端还有一个精细的发现：监测探针应该用\textbf{原始梯度}而不是
动量平滑后的更新方向——前者的顶部特征向量对齐显著更高
（$0.563\pm0.017$ vs $0.511\pm0.024$，CI 不重叠）。这把``监测
$\lmax$''与 CDAT/Kalra 式``调步长''在方法论上分叉开：前者要梯度、
后者要（预条件）更新方向。

\subsection{控制端：三种失败形态与一个例外}\label{sec:ch01-control}

第二条腿（第五至八章解读）是论文的主体。论文造了一系列尖锐度
控制器并观察它们的结局：

\begin{enumerate}
  \item \textbf{代理驱动的范数帽}（C1 权重范数帽、C2 $\gamma$ 帽）：
        投影确实频繁触发，$\lmax$ 纹丝不动——\textbf{帽原地不动}；
  \item \textbf{秩-1 plant}（C3，直接改写梯度、沿顶部特征向量阻尼）：
        结局取决于特征向量估计 $\hat v$ 的精度，两种体制都不下逃逸——
        估计滞后（体制 i）则周期-2 振荡残留、自稳定帽保持、钉在
        $2/\eta$；估计精确（体制 ii）则振荡被解除武装、帽被移除、
        $\lmax$ \textbf{上逃逸}冲过阈值（MLP 上逃到 $261.8\pm5.4$，
        阈值是 $40$）；
  \item \textbf{规模版秩-1 plant}（ResNet-18/CIFAR-100，修复感知伪
        收敛后）：执行器全力触发，$\lmax$ 与 SGD 统计不可分——
        \textbf{有动作仍被钉住}：尖锐度在秩-$(n{-}1)$ 正交补（与
        顶部特征向量垂直的全部参数方向构成的空间）上
        重新分布，单个方向压不住高秩谱。
\end{enumerate}

唯一的例外是 SAM：它的扰动
$\varepsilon=\rho\mathbf{g}/\|\mathbf{g}\|$（$\rho$ 为 SAM 的扰动
半径超参，决定在梯度方向上探出多远）作用在\textbf{每一个}
方向上，在所有设置中都能把 $\lmax$ 压到阈值之下（MLP 上
$21.35\pm0.04 < 40$；ResNet-18/CIFAR-100 上 $22.63\pm1.73$，与
SGD/秩-1 的 $\sim 29$ CI 分离）。

围绕``为什么秩-1 失败''，论文布置了四道机制排除防线：反驳``谱迁移''
替代解释的多尖锐景观测试（第七章解读）、正交阻尼与滞后拆分
（第七章）、排除钳位假象的 clamp-free 桥（第七章）、以及检验
X-SAM（第三方提出的``SAM $+$ 秩-1 顶部方向修正''组合优化器，Duan
等 2026，第六章详解）秩-1 修正贡献的忠实消融（第六章，$n{=}10$：其效应量上界
$+1.85$——效应量即真实差异的大小，上界意为``就算有效应也不超过
此值''——低于确定性噪声地板、无精度增益）。理论层面，第九章解读的
自稳定帽命题把这三种失败形态统一为一个二维递推的代数事实：
\textbf{不存在 $y^{*}\le 0$ 的翻号周期-2 轨道}（翻号 = 迭代每步改变
符号；周期-2 轨道 = 参数每两步一循环的序列，即 EoS 振荡的学名，
见第九章）——帽只能被顶开或
原地保留，永远压不下去。

\subsection{一张读懂全文的地图}\label{sec:ch01-map}

\begin{table}[!ht]
\centering\small
\caption{论文四项贡献与本解读章节的对应。每个贡献都由``实验证据 +
机制解释 + 适用域声明''三件套构成。}
\label{tab:ch01-map}
\begin{tabular}{p{0.46\linewidth}p{0.22\linewidth}p{0.22\linewidth}}
\toprule
贡献 & 论文章节 & 本解读章节 \\
\midrule
单 HVP 的 $\lmax$ 监测器与代理阶梯 & \S4 & 第四章 \\
监测探针是梯度而非更新方向 & \S4.3 & 第四章 \\
预测器不是执行器（C1--C3 + SAM + 四道压力测试） & \S5 & 第五--七章 \\
阈值不可局部下压、机制随体制变化 & \S6 理论 & 第九章 \\
规模复现（ResNet-18/CIFAR-100） & \S5.7 & 第八章 \\
实践建议与七条局限 & \S7 & 第十章 \\
\bottomrule
\end{tabular}
\end{table}

读图提示：第五至八章是一个递进式压力测试（escalating stress test）序列——每当你想对``秩-1 失败''提出一个替代解释（也许执行器
没触发？也许钳位锁住了它？也许只是 MLP 玩具设置的现象？），下一章
的实验就是为排除这个解释而设计的。这种``把审稿人的质疑提前做成
实验''的写法，是论文方法论上最值得学习的部分。

\subsection{为什么这篇论文值得细读}\label{sec:ch01-why}

三个理由。其一，\textbf{负结果有用}：它划清了 EoS 工程化利用的
边界——监测可行、单方向控制不可行——为后续优化器设计省下大量
试错成本；并且论文展示了失败是\textbf{结构性的}（命题 1 的代数
事实），不是调参没调好。其二，\textbf{方法学示范}：论文两次报告
自己的测量错误（幂迭代伪收敛低估 $45\%$、钳位门控制造休眠假象），
并设计了专门实验自我纠错——\texttt{FINDINGS\_zh.md} 日志保留了
完整的演化轨迹。其三，\textbf{反直觉的附属发现}：上逃逸到
$\lmax\approx 262$ 的臂在 held-out（留出）测试上反而是三臂最优
（MSE $0.0932$ vs SGD $0.1142$ vs SAM $0.1296$）——最尖的臂泛化
最好，逆着朴素平坦度-泛化信条（第十章局限 (vi) 解读）。

\subsection*{本章要点}
\begin{itemize}
  \item 论文把``EoS 可否工程化利用''切成预测/控制两半：预测端交付
        单 HVP 监测器（EoS 体制 $+0.981\pm0.030$），控制端证明
        单方向干预结构性失败。
  \item 控制失败有三种形态——帽原地不动、上逃逸、有动作仍被钉住
        ——共同点是永远压不到阈值以下；唯 SAM 全局动作下行。
  \item 证据链三层：控制实验、机制排除（迁移/滞后/钳位/规模）、
        2-D 约定理；每层适用域明确。
  \item 论文的价值不止结论：自我纠错的方法学与``最尖的臂泛化
        最好''的反直觉发现同样值得细读。
\end{itemize}
